\documentclass[../main.tex]{subfiles}
\begin{document}

\chapter{Păduri de mulțimi disjuncte}

Acest capitol nu conține teorie, deocamdată.

Fac o precizare legată de nomenclatură. Nu îmi plac denumirile larg răspîndite:

\begin{itemize}
  \item \textit{union-find}, pentru că nu îmi place să identific structura de date prin operații. Este ca și cum am denumi arborii de intervale \textit{query-update}.
  \item \textit{DSU} (\textit{disjoint set union}), din nou pentru că include numele unei operații.
\end{itemize}

În aceste note voi folosi denumirea DSF (engl. \textit{disjoint-set forest}), care consideră că arată mult mai clar ce reprezintă structura: o colecție arbori disjuncți care partiționează o colecție de obiecte.

Urmează cîteva probleme de care m-am izbit de-a lungul anilor și care folosesc DSF în moduri mai puțin evidente. Vă recomand și probleme din alte capitole care folosesc DSF:

\begin{itemize}
  \item \hyperref[problem:extending-a-set-of-points]{Extending a Set of Points}
  \item \hyperref[problem:mst-for-each-edge]{Minimum Spanning Tree for Each Edge}
  \item \hyperref[problem:regate]{Regate}
  \item \hyperref[problem:0-1-mst]{0-1 MST}
\end{itemize}

\end{document}
